\chapter*{Resumen}

\section*{\tituloPortadaVal}

Hoy en día, muchos programas de televisión, sobre todo los informativos, utilizan elementos de realidad aumentada (RA) que se muestran en pantalla para complementar la información ofrecida por los locutores. Gracias a esta tecnología es posible integrar gráficos en tres dimensiones, datos interactivos o visualizaciones complejas directamente en el plató, de manera que parecen formar parte del espacio real. Con ello no solo se completa y facilita la comprensión de los contenidos, sino que también se aporta un aspecto más dinámico y atractivo a la emisión.

En este Trabajo de Fin de Grado se estudia cómo aplicar la realidad aumentada a las retransmisiones en directo mediante plataformas como YouTube u otros servicios de \textit{streaming}. Para llevarlo a cabo se emplea OBS Studio como herramienta principal de producción, analizando de qué forma puede adaptarse para generar efectos similares a los que se utilizan en televisión profesional.

El trabajo se centra especialmente en el caso de concursos de programación en línea, donde la RA puede ser de gran ayuda para presentar la información de manera más clara, facilitar la interacción con el público y mejorar el valor visual del contenido. De esta forma se busca demostrar cómo esta tecnología puede enriquecer tanto la experiencia del espectador como la del propio presentador.

\section*{Palabras clave}

\noindent Realidad Aumentada, Retransmisión en directo, complemento OBS, Modelos 3D, Marcadores ArUco, OpenCV, DOMjudge


